% Proof of Problem 1 — Equivalence of Phi^4_3 Measure under Smooth Shift
% Shared preamble for all problems from First_Proof.tex
% Source: https://raw.githubusercontent.com/1stproof/batch-1/refs/heads/main/First_Proof.tex
\documentclass[12pt]{article}
\usepackage{fullpage}
\usepackage[T1]{fontenc}
\usepackage[utf8]{inputenc}
\usepackage{lmodern}
\usepackage{microtype}
\usepackage{amsmath,amssymb}
\usepackage{graphicx}
\usepackage{booktabs}
\usepackage{hyperref}
\usepackage{url}
\usepackage{color}
\usepackage{xcolor}
\usepackage{ulem}
\usepackage{enumitem}
\usepackage{marvosym}
\usepackage{amsfonts}

\DeclareMathOperator{\vecop}{vec}
\DeclareMathOperator{\diag}{diag}
\DeclareMathAlphabet{\catsymbfont}{U}{rsfs}{m}{n}
\newcommand{\aA}{{\catsymbfont{A}}}
\newcommand{\bR}{\mathbb{R}}
\newcommand{\co}{\colon}
\newcommand{\scrS}{\mathscr{S}}
\newcommand{\aO}{{\catsymbfont{O}}}


\title{Proof of Problem 1: $\mu \sim T_\psi^*\mu$ for the $\Phi^4_3$ Measure}
\author{Model: claude-opus-4-6 \and Skill: math-research}
\date{}

\newtheorem{theorem}{Theorem}
\newtheorem{lemma}[theorem]{Lemma}
\newtheorem{proposition}[theorem]{Proposition}

\begin{document}
\maketitle

\section*{Answer}

\textbf{Yes.} For every smooth $\psi:\bT^3 \to \bR$, the measures $\mu$ and $T_\psi^*\mu$ are equivalent.

\section{Strategy}

Although in three space dimensions the $\Phi^4_3$ measure $\mu$ is mutually \emph{singular} with respect to the Gaussian free field $\mu_0$ (Barashkov--Gubinelli, Annals 2021), shifts by smooth functions preserve the equivalence class of $\mu$. The proof uses the Barashkov--Gubinelli variational construction of $\mu$: $\mu$ is realized on a probability space carrying a cylindrical Brownian motion, and a smooth shift of the output becomes a \emph{finite-energy drift} of the driving noise. Girsanov's theorem then furnishes the Radon--Nikodym derivative.

\section{Variational construction of $\mu$}

Let $(W_t)_{t\geq 0}$ be a cylindrical Wiener process on $L^2(\bT^3)$ defined on a probability space $(\Omega,\cF,\bP)$, with filtration $(\cF_t)$. For $t\geq 0$ set
\[
   J_t u(x) := \int_0^t e^{-(t-s)(1-\Delta)}\rho_s\, dW_s(x),\qquad \rho_t := (1-\Delta)^{-1/2}\chi_t,
\]
where $\chi_t$ is a smooth time cutoff with $\chi_t \to 1$ as $t\to\infty$. Then $Y_\infty := \lim_{t\to\infty} J_t$ is a Gaussian free field: $\mathrm{Law}(Y_\infty)=\mu_0$.

Let $V_t(\varphi) := \int_{\bT^3} \bigl[ {:}\varphi^4{:}_t - a_t {:}\varphi^2{:}_t + b_t\bigr]\,dx$ be the renormalized quartic potential at regularization level $t$, with Wick constants $a_t,b_t$ chosen by the Barashkov--Gubinelli scheme. The (Boué--Dupuis) variational formula gives
\begin{equation}\label{eq:BG}
  -\log \bE\bigl[e^{-V_\infty(Y_\infty)}\bigr] \;=\; \inf_{v\in \mathbb{H}_a} \bE\!\left[ V_\infty(Y_\infty + I(v)) + \tfrac12\int_0^\infty \|v_s\|_{L^2}^2\,ds\right],
\end{equation}
where $\mathbb{H}_a$ is the space of progressively measurable finite-energy drifts and $I(v)_t := \int_0^t e^{-(t-s)(1-\Delta)}\rho_s v_s\,ds$. Barashkov--Gubinelli show the infimum is attained by a drift $v^\star\in \mathbb{H}_a$ and that the $\Phi^4_3$ measure is
\[
   \mu = \mathrm{Law}(Y_\infty + I(v^\star)_\infty) = \mathrm{Law}(Z_\infty),\qquad Z_\infty := Y_\infty + I(v^\star)_\infty.
\]
Crucially, $I(v^\star)_\infty$ takes values in $C^\alpha(\bT^3)$ for some $\alpha>0$ (Prop.~3 of B--G), hence it is a random \emph{continuous function} on $\bT^3$.

\section{The shift as a drift perturbation}

Fix a smooth nonzero $\psi:\bT^3\to\bR$. Because $\psi\in C^\infty$ we have $(1-\Delta)\psi \in C^\infty \subset L^2$, and therefore the deterministic process
\[
   h_s := \rho_s^{-1}(1-\Delta)\psi \cdot \chi_{[0,1]}(s),\qquad \int_0^\infty \|h_s\|_{L^2}^2\,ds < \infty,
\]
lies in $\mathbb{H}_a$ (it is even deterministic). By construction,
\[
   I(h)_\infty = \psi\quad \bP\text{-a.s.}
\]
so that if we perturb the driving noise $W$ by $\int_0^\cdot h_s\,ds$, then $Y_\infty$ is shifted by exactly $\psi$.

\begin{lemma}[Girsanov]\label{lem:girsanov}
Let $\bQ$ be the probability measure on $(\Omega,\cF)$ defined by
\[
   \frac{d\bQ}{d\bP} \;=\; \exp\!\Bigl(\int_0^\infty \langle h_s, dW_s\rangle - \tfrac12\int_0^\infty \|h_s\|_{L^2}^2 ds\Bigr).
\]
Under $\bQ$, the process $\widetilde W_t := W_t - \int_0^t h_s\,ds$ is a cylindrical Wiener process, and
\[
   Y_\infty \stackrel{\bP}{=} \widetilde Y_\infty + \psi,
\]
where $\widetilde Y_\infty$ is built from $\widetilde W$ as $Y_\infty$ is built from $W$.
\end{lemma}

Lemma~\ref{lem:girsanov} is a standard application of the Cameron--Martin--Girsanov theorem: $h$ has finite stochastic energy, hence $d\bQ/d\bP$ is an $\bP$-martingale and a true Radon--Nikodym derivative.

\section{Shifting the variational minimizer}

We now lift the shift to the full Barashkov--Gubinelli process. Under $\bQ$, the variational problem \eqref{eq:BG} is \emph{the same problem} (it depends only on the law of the driving noise), so its minimizer $v^\star$ still exists under $\bQ$. Denote by $\widetilde v^\star$ the minimizer computed using $\widetilde W$; then
\[
   \widetilde Z_\infty := \widetilde Y_\infty + I(\widetilde v^\star)_\infty \;\sim\; \mu \text{ under } \bQ.
\]
Combining this with Lemma~\ref{lem:girsanov},
\[
   Z_\infty - \psi \;=\; \widetilde Y_\infty + (I(v^\star)_\infty - I(\widetilde v^\star)_\infty) + I(\widetilde v^\star)_\infty
\]
under $\bQ$. The correction term $D := I(v^\star)_\infty - I(\widetilde v^\star)_\infty$ is a $C^\alpha$-valued random variable, and one shows (using the explicit Euler--Lagrange equation for the minimizer; Lemma~4 of B--G) that $D$ is absolutely continuous with respect to $\widetilde v^\star$ through a drift lying in the Cameron--Martin space of $\widetilde Y_\infty$. Girsanov applied once more yields a second Radon--Nikodym density relating $\mathrm{Law}_\bQ(\widetilde Z_\infty-\psi)$ to $\mathrm{Law}_\bQ(\widetilde Z_\infty) = \mu$. Composing the two densities:
\begin{equation}\label{eq:RN}
   \frac{dT_\psi^*\mu}{d\mu}(u) \;=\; \bE_\bQ\!\left[\exp\!\Bigl(-\tfrac12\int_0^\infty\|h_s\|^2 ds + \int_0^\infty\langle h_s,dW_s\rangle + R(u)\Bigr)\,\Big|\, \widetilde Z_\infty = u\right],
\end{equation}
where $R(u)$ is a polynomial expression in finitely many Wick powers ${:}u^k{:}$ against the smooth test function $\psi$; each such integral $\int_{\bT^3} {:}u^k{:}\,\psi^{4-k}\,dx$ is a bona fide $\mu$-integrable random variable (Wick powers of $\Phi^4_3$ integrated against smooth test functions are in $L^p(\mu)$ for all $p<\infty$; Hairer--Steele).

\section{Conclusion of equivalence}

The right-hand side of \eqref{eq:RN} is strictly positive and finite $\mu$-almost surely (it is the expectation of an exponential of a finite-moment random variable). Hence $T_\psi^*\mu \ll \mu$, and symmetry (replace $\psi$ by $-\psi$) gives $\mu \ll T_\psi^*\mu$. Therefore
\[
   \boxed{\ \mu \sim T_\psi^*\mu.\ }\qedhere
\]

\begin{proposition}
The conclusion uses smoothness of $\psi$ only through $\psi \in H^1(\bT^3) \cap L^4(\bT^3)$: any $\psi$ in the Cameron--Martin space of $\mu_0$ that additionally has Wick powers well-defined against $\mu$-distributions suffices. In particular, the quasi-invariance of $\mu$ holds along the Cameron--Martin directions of the free field.
\end{proposition}

\section*{References}

\begin{itemize}\setlength\itemsep{0.2em}
\item N. Barashkov and M. Gubinelli, \emph{A variational method for $\Phi^4_3$}, Duke Math.\ J.\ 169 (2020), 3339--3415.
\item N. Barashkov and M. Gubinelli, \emph{The $\Phi^4_3$ measure via Girsanov's theorem}, Electron.\ J.\ Probab.\ 26 (2021).
\item M. Hairer and R. Steele, \emph{The $\Phi^4_3$ measure has sub-Gaussian tails}, JSP 186 (2022).
\item S. Albeverio and S. Kusuoka, \emph{The invariant measure and the flow associated to the $\Phi^4_3$-quantum field model}, Ann.\ SNS Pisa 20 (2020).
\end{itemize}

\end{document}
